% =============================================================================
%  Parameter selection and error analysis.
%  SHARED: \input by both paper.tex (as Section 5) and technical-report.tex,
%  so the two documents cannot report different numbers.
% =============================================================================
\section{Parameter Selection: From Diagnosis to a Feasible Operating Point}
\label{sec:params}

The FLPSI layer is parameterised by the sub-sample width $w$, the number of sub-samples $T$, and
the acceptance threshold $t$ on the number of colliding sub-samples (earlier sections write
$\theta$ for the identical quantity; we use $t$ here to match the hypothesis-testing language of
\S\ref{sec:design}). This section shows that $(w,t,T)$ cannot be chosen from an assumed impostor
model --- doing so at the naïve Super-Bit code understates the false-accept rate by three orders
of magnitude and hits a floor no threshold can cross --- but that this is a fact about \emph{that
code}, not about the protocol. Once the code is co-designed with the accept rule, exactly the same
arithmetic yields a feasible, verified operating point.

\subsection{The accept model}
\label{sec:accept-model}

A sub-sample selects $w$ of the $d = 128$ positions uniformly at random and collides when the two
codes agree on all of them. Writing $H = \|\xt \oplus \yt_i\|_1$ for the Hamming distance between
probe and enrolled code, the per-sub-sample collision probability and the acceptance probability
over $T$ independent sub-samples are
\begin{equation}
\label{eq:accept}
q(H) \;=\; \frac{\binom{d-H}{w}}{\binom{d}{w}} \;\approx\; \Bigl(1-\tfrac{H}{d}\Bigr)^{w},
\qquad
\Pr[\text{accept} \mid H] \;=\; \Pr\bigl[\mathrm{Binomial}(T, q(H)) \ge t\bigr].
\end{equation}
Acceptance depends on a pair \emph{only through $H$}, so the false-accept and false-reject rates
are exact functionals of the two Hamming histograms $P_{\mathrm{imp}}, P_{\mathrm{gen}}$,
\begin{equation}
\label{eq:rates}
\mathrm{FAR} = \sum_{H} P_{\mathrm{imp}}(H)\,\Pr[\text{accept}\mid H],
\qquad
\mathrm{FRR} = 1 - \sum_{H} P_{\mathrm{gen}}(H)\,\Pr[\text{accept}\mid H].
\end{equation}
As shown in \S\ref{sec:design}, thresholding the sub-sample count $K$ is not an arbitrary choice:
it is the Neyman--Pearson-optimal test for the sufficient statistic FLPSI already computes, given
the two aggregate per-trial accept masses $q_g = \mathbb{E}_{P_{\mathrm{gen}}}[q(H)]$ and
$q_i = \mathbb{E}_{P_{\mathrm{imp}}}[q(H)]$. That result licenses treating parameter selection as a
one-dimensional problem --- choosing $T$ and $t$ once $q_g, q_i$ are known --- and the Gaussian
form derived there,
\begin{equation}
\label{eq:tmin}
T_{\min} \;\approx\; \frac{\bigl[\,z_\alpha\sqrt{q_i(1-q_i)} + z_\beta\sqrt{q_g(1-q_g)}\,\bigr]^2}
                            {(q_g - q_i)^2},
\end{equation}
is the design rule used throughout this section: it says $T$ is governed entirely by the gap
$q_g - q_i$, which in turn is governed entirely by the impostor and genuine Hamming laws
(\S\ref{sec:design}, Lemma~2.1 / Prop.~3.2). Every number below plugs the \emph{measured} laws
into~\eqref{eq:accept}--\eqref{eq:tmin}; nothing is assumed.

\subsection{Diagnosis: what an assumed impostor law costs}
\label{sec:diagnosis}

The naïve design fixes $(w,t) = (14,2)$ under the assumption that impostor codes are uniform,
$H \sim \mathrm{Binomial}(128, \tfrac12)$, i.e.\ mean $64$ and $\sigma = 5.66$. The Super-Bit
bridge (\S\ref{sec:design}) does center the mean correctly --- measured impostor mean $64.05$ ---
but its measured standard deviation is $\sigma = 20.11$, $3.6\times$ the assumed value, at
\emph{the same mean}. No first-moment check on this code would ever catch the problem.

This is exactly the pathology Lemma~2.1 (\S\ref{sec:design}) predicts: $q(H)$ is convex, so any
mean-preserving spread of the impostor law strictly increases $q_i$, and hence $T_{\min}$, and
hence communication. Concretely, evaluating~\eqref{eq:rates} at $(w,t)=(14,2)$ against the assumed
uniform law reproduces the design-time false-accept rate of $2.954\times10^{-5}$ exactly; against
the \emph{measured} Super-Bit law the same expression gives $4.578\times10^{-2}$ --- a
\textbf{1,550$\times$ gap} between the assumed and the actual per-record error, entirely explained
by the excess variance.

\paragraph{The collision floor.} Since $q(0)=1$, every sub-sample collides whenever $H=0$, so
$\Pr[\text{accept}\mid H{=}0]=1$ for \emph{every} $(w,t,T)$: a code with any zero-Hamming impostor
pairs has a floor on per-record FAR that no parameter choice can go below. On the Super-Bit code,
$2$ of $1{,}438{,}800$ held-out probe$\times$gallery impostor pairs collide exactly, giving
\begin{equation}
\label{eq:floor}
p_0 = 1.39\times10^{-6}, \qquad \text{exact Poisson 95\% CI } [1.68\times10^{-7},\, 5.02\times10^{-6}],
\end{equation}
which alone contributes $\mathrm{FAR}_{\text{query}} = 1-(1-p_0)^N \approx 6.9\times10^{-3}$ at
$N=5{,}000$.\footnote{An earlier internal accounting reported $8$ collisions out of $2{,}877{,}600$
pairs ($2.78\times10^{-6}$); that count additionally includes enrolment-versus-enrolment pairs,
which are never issued as queries. We retire that figure and use the probe$\times$gallery
definition~\eqref{eq:floor} throughout, as the operationally meaningful one.} By
Corollary~6.1 (\S\ref{sec:design}), $p_0$ is a coding question, not a protocol question: it is
governed by how many pairs of identities the code's \emph{effective} dimension $m_{\mathrm{eff}}$
can keep apart, not by the nominal length $d$. This is why raising $d$ does not help on its own ---
the dimension sweep shows Super-Bit's $\sigma$ barely moves ($20.1 \to 18.1$) as $D:16\to64$ ---
and it is the reason this section does not end with a retuning sweep, but with a different code.

\textbf{We treat both numbers above as a diagnosis, not a verdict.} They say precisely what a code
must fix --- shrink the impostor variance toward the uniform floor $d/4$ --- and \S\ref{sec:fix}
shows a code that does.

\subsection{Effective bits track the impostor variance}
\label{sec:effbits}

The Daugman binomial degrees-of-freedom estimator translates a measured impostor
$\sigma$ into an effective code length: at balance ($p=\tfrac12$), $m_{\mathrm{eff}} =
d^2/(4\sigma^2)$~\cite{Daugman03}, which recovers $m_{\mathrm{eff}} = d$ exactly for the ideal
uniform code ($\sigma=\sqrt{d}/2$). Table~\ref{tab:effbits} applies it to the three codes measured
in this paper.

\begin{table}[htbp]
\centering
\caption{Impostor standard deviation and Daugman effective bits, $d=128$. The uniform-ideal
reference is $\sigma=5.66$, $m_{\mathrm{eff}}=128$.}
\label{tab:effbits}
\small
\begin{tabular}{@{}lrrl@{}}
\toprule
code & impostor $\sigma$ & $m_{\mathrm{eff}} = d^2/(4\sigma^2)$ & agreement with an independent estimator \\
\midrule
Super-Bit (naïve)      & 20.11 & $\approx 10$   & participation ratio $9.84$ (measured) \\
R1: metric-aware bridge & 14.60 & $\approx 18$   & participation ratio $13.6$ (measured) \\
R2: feature-head code   & 6.72  & 90.7 & participation ratio $84.6$ (measured) \\
\bottomrule
\end{tabular}
\end{table}

Super-Bit's $m_{\mathrm{eff}}\approx10$ is not a defect of the bridge but a consequence of coding a
16-dimensional embedding into $128$ correlated bits; the bridge is faithful to the embedding's
angular structure, and that structure is what has low effective dimension. A distribution-free
alternative to the Daugman estimate is the coding-rate functional of maximal coding-rate reduction
(MCR$^2$)~\cite{YuEtAl20}, $R(Z,\varepsilon) = \tfrac12\log\det\bigl(I + \tfrac{p}{n\varepsilon^2}
ZZ^\top\bigr)$, which counts effectively independent directions without assuming a binomial fit;
quantifying $R(Z,\varepsilon)$ and the rate reduction $\Delta R$ for the three codes on the
held-out embeddings directly is left to future work.

\subsection{The fix: co-designing the code raises the effective bits and closes the gap}
\label{sec:fix}

Corollary~4.2 and Theorem~4.3 (\S\ref{sec:design}) turn \S\ref{sec:diagnosis}'s diagnosis into a
training objective: balance every bit ($\mathbb{E}[D_i]=\tfrac12$) and decorrelate bit pairs
($\mathrm{Cov}(D_i,D_j)\to0$), which is exactly minimizing the impostor variance the accept model
is sensitive to. Applied to the same 128-bit FLPSI backend, this yields two successive codes whose
impostor $\sigma$ moves toward the uniform ideal of $5.66$ and whose measured error moves with it:

\begin{table}[htbp]
\centering
\caption{Impostor $\sigma$, EER, and zero-Hamming collisions across the diagnosis-to-fix ladder
(SOCOFing, Easy probes, 1,200 held-out identities, identity-bootstrap 95\% CI).}
\label{tab:ladder}
\small
\begin{tabular}{@{}lrrr@{}}
\toprule
code & impostor $\sigma$ & EER (95\% CI) & $H{=}0$ collisions \\
\midrule
Super-Bit (naïve)       & 20.11 & 1.89\% [1.57, 2.22] & 2 / 1{,}438{,}800 \\
R1: metric-aware bridge & 14.60 & 0.87\% [0.62, 0.99] & 0 \\
R2: feature-head code   & 6.72  & 0.17\% [0.06, 0.26] & 0 \\
\bottomrule
\end{tabular}
\end{table}

Both co-designed codes drive the collision floor of \S\ref{sec:diagnosis} to zero on the held-out
set --- the same diagnostic that indicted Super-Bit now certifies the fix.

\paragraph{A feasible operating point.} Plugging the feature-head code's measured $q_g, q_i$ into
the $T_{\min}$ rule~\eqref{eq:tmin} and rounding to the nearest integer, feasible parameters at a
target TAR $\geq 95\%$ and per-record FAR $\leq 10^{-2}$ are $(w,t,T) = (6,3,18)$ --- an order of
magnitude fewer sub-samples than Super-Bit's published $T=64$, because the tighter impostor law
raises $q_g - q_i$ and shrinks the numerator of~\eqref{eq:tmin} at once. Equation~\eqref{eq:accept}
predicts a per-record FAR of $8.12\times10^{-3}$ at this point; running the \emph{real}
\texttt{flash-psi} binary (masked-OPRF + garbled circuit + VOLE + Shamir sharing, not a
closed-form simulation) at $(6,3,18)$ gives real TAR $97.0\%$, real FAR $0.80\%$
($8.02\times10^{-3}$) against the $8.12\times10^{-3}$ prediction, $\sim\!84$\,ms/query.
This finite-$T$ agreement between
the asymptotic design rule and the deployed binary is the same pattern observed at Super-Bit's own
published point (\S\ref{sec:design}, finite-$T$ caveat), where real and predicted FAR matched to
the digit; here it certifies that the design rule transfers to the co-designed code, not just the
naïve one.

\subsection{Communication follows the same variance that fixed accuracy}
\label{sec:comms}

Because online communication scales with $T$ (\S\ref{sec:design}, assumption A4), the same
variance reduction that lowered EER lowers bytes. Against Super-Bit's original, uncalibrated
$T=64$ point, online communication measured on the real binary is $7953$\,KB per query at
$N=5{,}000$ --- and Super-Bit has no feasible \emph{matched} operating point at all, since no
$(w,t,T)$ reaches per-record FAR $\leq 10^{-2}$ on its code, so no like-for-like comparison to it
exists. Between the two codes that \emph{do} have a matched point (TAR $\geq 95\%$, per-record
FAR $\leq 10^{-2}$, $N=5{,}000$), R1 metric-aware bridge at $(w,t,T)=(19,3,34)$ costs
$4317$\,KB while R2 feature-head reaches the same target at $T=18$
(i.e.\ $(w,t,T)=(6,3,18)$) for $2377$\,KB --- a $45\%$ reduction at the same target
operating point. These two
comparisons must never be conflated: the $7953\to2377$\,KB figure is against an untuned $T=64$
baseline that has no matched point of its own, while the $-45\%$ figure is the like-for-like
comparison between two calibrated codes on the same protocol. The latter is the paper's
communication contribution; the former only shows how far an uncalibrated parameterisation is from
either.

\paragraph{Summary.} Table~\ref{tab:ladder} and the operating point $(6,3,18)$ are the same
object: parameters are not chosen against an assumed law and then defended, they are read off the
measured impostor Hamming distribution of whichever code is deployed, and a better code is a
strictly better choice of $(w,t,T)$ for free.
